\mxcorrections{7}\label{C07-CORR}
\textbf{Usage professeur.} Les réponses qui suivent concernent les tâches non résolues dans le cours. Faire d’abord expliciter la démarche ; fournir une aide graduée avant la solution. Les durées sont des estimations à étalonner en classe.
\subsection*{Diagnostic et remédiation}
\begin{enumerate}
\item $0$. \textit{Action :} Reprendre le critère d’orthogonalité.
\item Oui : $2-2=0$. \textit{Action :} Substituer x et y simultanément.
\item $x=3$, $y=2$. \textit{Action :} Additionner puis remplacer.
\item $5$. \textit{Action :} Calculer le vecteur puis sa norme.
\item $(x-2)^2-4$. \textit{Action :} Développer pour vérifier.
\item Le seul couple $\pt{0;0}$. \textit{Action :} Utiliser la somme de deux carrés positifs.
\end{enumerate}
\subsection*{Activité de recherche}
Les essais suggèrent H(2;4), mais Pythagore dans PHM prouve la minimalité. Le projeté se calcule avec la normale (2;-1).
\subsection*{Applications immédiates}
\begin{enumerate}
\item $3(x-2)+4(y+1)=0$, soit $3x+4y-2=0$.
\item $H\pt{2;1}$ et distance 2 ; la normale est horizontale.
\item $(x+1)^2+(y-4)^2=0$ : le seul point $\pt{-1;4}$, pas un cercle de rayon strictement positif.
\item $(x-1)^2=0$, donc seul point $\pt{1;1}$ ; la droite est tangente au cercle.
\end{enumerate}
\subsection*{Exercices : démarche et résultats}
\paragraph{Exercice 7.1 -- Construire une droite.} $2(x+1)+5(y-3)=0$, soit $2x+5y-13=0$. Directeur $\pt{5;-2}$, donc $B\pt{4;1}$ est sur la droite ; contrôle $8+5-13=0$.
\paragraph{Exercice 7.2 -- Calculer un projeté.} Normale $\pt{2;-1}$ ; $\lambda =\dfrac{8-3}{5}=1$. $H=\pt{4;3}-\pt{2;-1}=\pt{2;4}$. $2(2)-4=0$ et $\vect{PH}=\pt{-2;1}$ est normal ; distance $\sqrt5$.
\paragraph{Exercice 7.3 -- Reconnaître un cercle.} $(x-3)^2+(y+2)^2=25$ ; centre $\pt{3;-2}$, rayon 5. Pour $y=-2$, $x=3\pm5$ : $\pt{-2;-2}$ et $\pt{8;-2}$.
\paragraph{Exercice 7.4 -- Médiatrice.} Milieu I(1;2), normale $\vect{AB}=\pt{4;2}$. $4(x-1)+2(y-2)=0$, soit $2x+y-4=0$. I vérifie $2+2-4=0$.
\paragraph{Exercice 7.5 -- Le meilleur accès.} Le projeté H(2;4) minimise PH, de longueur $\sqrt5$ km. Pour tout M sur la route, le triangle PHM est rectangle en H et $PM^2=PH^2+HM^2\geq PH^2$. Égalité pour M=H.
\paragraph{Exercice 7.6 -- Paramètre et ensemble.} $(x-1)^2+(y+2)^2=5-k$. Si $k<5$, cercle de centre $\pt{1;-2}$ et rayon $\sqrt{5-k}$ ; si $k=5$, point ; si $k>5$, vide.
\paragraph{Exercice 7.7 -- Une parabole comme lieu.} $MF^2=x^2+(y-2)^2$, distance à l’axe au carré $y^2$. Égalité : $x^2-4y+4=0$, soit $y=\dfrac{x^2}{4}+1$. Les distances positives permettent la réciproque.
\paragraph{Exercice 7.8 -- Reprise différée.} Centre $\pt{2;0}$, rayon 2 ; $(x-2)^2+y^2=4$. Pour x=5, $y^2=-5$, impossible. Interprétation : droite à distance 3 du centre, supérieure au rayon.
\paragraph{CQFD, exercice 18 p. 279, questions a, b, c et d.} Des réponses possibles sont $\pt{-1;2}$, $\pt{3;-6}$, $\pt{\dfrac14;-1}$ et $\pt{1;0}$. Tout multiple non nul convient. Pour $x=4$, écrire $x-4=0$.
\paragraph{CQFD, exercice 21 p. 279, questions a et b.} a : $x-(y-2)=0$, soit $x-y+2=0$. b : $4(x+1)+6(y-3)=0$, soit $2x+3y-7=0$. Substituer les coordonnées des points confirme l’appartenance.
\subsection*{QCM : réponses et raisons}
\textbf{QCM 1.}
\begin{enumerate}
\item B. Coefficients de x et y.
\item A. Produit avec la normale nul.
\item D. Lire les coordonnées opposées aux signes.
\item C. Racine positive de 9.
\item B. Somme de carrés non négative.
\end{enumerate}
\textbf{QCM 2.}
\begin{enumerate}
\item C. Même abscisse, ordonnée nulle.
\item A. Écart vertical absolu.
\item D. Normale non nulle.
\item B. Les deux carrés doivent être nuls.
\item C. $y^2=0$.
\end{enumerate}
\subsection*{Vrai ou faux}
\begin{enumerate}
\item Vrai : par P passe une unique perpendiculaire à la droite ; si P appartient à la droite, H=P.
\item Faux : si c=0, un point ; si c>0, ensemble vide.
\item Faux : la substitution d’un paramétrage de droite dans l’équation de cercle donne une équation quadratique de coefficient dominant strictement positif, donc au plus deux racines.
\end{enumerate}
\paragraph{Différenciation.} Pour un blocage technique, proposer la première étape seulement (domaine, facteur ou univers de référence), puis retirer l’aide. Pour un élève autonome, demander une seconde méthode et le traitement d’un cas limite. À l’oral, autoriser une reformulation initiale dans la langue maîtrisée, puis exiger une phrase mathématique précise en français.
